
\subsubsection{2.1 Communication Accommodation Theory}

Communication Accommodation Theory (CAT), introduced by Giles (1973), posits that individuals adjust their communication style to signal social identity and regulate social distance. Convergence---adapting toward an interlocutor's style---signals affiliation.

Niederhoffer and Pennebaker (2002) operationalized accommodation through Linguistic Style Matching (LSM), finding correlations around $r\approx 0.3$ between style matching and relationship quality in human-human dyads. This work established that accommodation is measurable through linguistic features and predicts interaction outcomes.

However, human-human accommodation research assumes symmetric agency. In human-AI interaction, this symmetry may not hold. AI systems accommodate through training objectives, not social motivation; humans may not perceive AI as warranting reciprocal accommodation.

\subsubsection{2.2 Human-AI Linguistic Adaptation}

Chen et al. (2026) analyzed 1,319 GPT-4o conversations from WildChat, demonstrating bidirectional linguistic accommodation with distinct temporal dynamics: model accommodation is front-loaded (strongest in early turns), while user convergence develops gradually.

This finding validates that accommodation exists in human-AI interaction but leaves the outcome relationship untested. The present work extends this line by testing whether accommodation magnitude predicts conversation continuation.

\subsubsection{2.3 Chatbot Engagement and the Uncanny Valley}

Ciechanowski et al. (2019) documented uncanny valley effects in chatbot interaction: users respond negatively to AI behaviors that approximate but do not fully achieve human-like communication. This suggests that excessive accommodation might have unintended consequences.

The inverted-U accommodation-engagement relationship found in the present study is consistent with this framework, though other explanations remain possible.

